% Options for packages loaded elsewhere
% Options for packages loaded elsewhere
\PassOptionsToPackage{unicode}{hyperref}
\PassOptionsToPackage{hyphens}{url}
\PassOptionsToPackage{dvipsnames,svgnames,x11names}{xcolor}
%
\documentclass[
  letterpaper,
  DIV=11,
  numbers=noendperiod]{scrartcl}
\usepackage{xcolor}
\usepackage{amsmath,amssymb}
\setcounter{secnumdepth}{5}
\usepackage{iftex}
\ifPDFTeX
  \usepackage[T1]{fontenc}
  \usepackage[utf8]{inputenc}
  \usepackage{textcomp} % provide euro and other symbols
\else % if luatex or xetex
  \usepackage{unicode-math} % this also loads fontspec
  \defaultfontfeatures{Scale=MatchLowercase}
  \defaultfontfeatures[\rmfamily]{Ligatures=TeX,Scale=1}
\fi
\usepackage{lmodern}
\ifPDFTeX\else
  % xetex/luatex font selection
\fi
% Use upquote if available, for straight quotes in verbatim environments
\IfFileExists{upquote.sty}{\usepackage{upquote}}{}
\IfFileExists{microtype.sty}{% use microtype if available
  \usepackage[]{microtype}
  \UseMicrotypeSet[protrusion]{basicmath} % disable protrusion for tt fonts
}{}
\makeatletter
\@ifundefined{KOMAClassName}{% if non-KOMA class
  \IfFileExists{parskip.sty}{%
    \usepackage{parskip}
  }{% else
    \setlength{\parindent}{0pt}
    \setlength{\parskip}{6pt plus 2pt minus 1pt}}
}{% if KOMA class
  \KOMAoptions{parskip=half}}
\makeatother
% Make \paragraph and \subparagraph free-standing
\makeatletter
\ifx\paragraph\undefined\else
  \let\oldparagraph\paragraph
  \renewcommand{\paragraph}{
    \@ifstar
      \xxxParagraphStar
      \xxxParagraphNoStar
  }
  \newcommand{\xxxParagraphStar}[1]{\oldparagraph*{#1}\mbox{}}
  \newcommand{\xxxParagraphNoStar}[1]{\oldparagraph{#1}\mbox{}}
\fi
\ifx\subparagraph\undefined\else
  \let\oldsubparagraph\subparagraph
  \renewcommand{\subparagraph}{
    \@ifstar
      \xxxSubParagraphStar
      \xxxSubParagraphNoStar
  }
  \newcommand{\xxxSubParagraphStar}[1]{\oldsubparagraph*{#1}\mbox{}}
  \newcommand{\xxxSubParagraphNoStar}[1]{\oldsubparagraph{#1}\mbox{}}
\fi
\makeatother

\usepackage{color}
\usepackage{fancyvrb}
\newcommand{\VerbBar}{|}
\newcommand{\VERB}{\Verb[commandchars=\\\{\}]}
\DefineVerbatimEnvironment{Highlighting}{Verbatim}{commandchars=\\\{\}}
% Add ',fontsize=\small' for more characters per line
\usepackage{framed}
\definecolor{shadecolor}{RGB}{241,243,245}
\newenvironment{Shaded}{\begin{snugshade}}{\end{snugshade}}
\newcommand{\AlertTok}[1]{\textcolor[rgb]{0.68,0.00,0.00}{#1}}
\newcommand{\AnnotationTok}[1]{\textcolor[rgb]{0.37,0.37,0.37}{#1}}
\newcommand{\AttributeTok}[1]{\textcolor[rgb]{0.40,0.45,0.13}{#1}}
\newcommand{\BaseNTok}[1]{\textcolor[rgb]{0.68,0.00,0.00}{#1}}
\newcommand{\BuiltInTok}[1]{\textcolor[rgb]{0.00,0.23,0.31}{#1}}
\newcommand{\CharTok}[1]{\textcolor[rgb]{0.13,0.47,0.30}{#1}}
\newcommand{\CommentTok}[1]{\textcolor[rgb]{0.37,0.37,0.37}{#1}}
\newcommand{\CommentVarTok}[1]{\textcolor[rgb]{0.37,0.37,0.37}{\textit{#1}}}
\newcommand{\ConstantTok}[1]{\textcolor[rgb]{0.56,0.35,0.01}{#1}}
\newcommand{\ControlFlowTok}[1]{\textcolor[rgb]{0.00,0.23,0.31}{\textbf{#1}}}
\newcommand{\DataTypeTok}[1]{\textcolor[rgb]{0.68,0.00,0.00}{#1}}
\newcommand{\DecValTok}[1]{\textcolor[rgb]{0.68,0.00,0.00}{#1}}
\newcommand{\DocumentationTok}[1]{\textcolor[rgb]{0.37,0.37,0.37}{\textit{#1}}}
\newcommand{\ErrorTok}[1]{\textcolor[rgb]{0.68,0.00,0.00}{#1}}
\newcommand{\ExtensionTok}[1]{\textcolor[rgb]{0.00,0.23,0.31}{#1}}
\newcommand{\FloatTok}[1]{\textcolor[rgb]{0.68,0.00,0.00}{#1}}
\newcommand{\FunctionTok}[1]{\textcolor[rgb]{0.28,0.35,0.67}{#1}}
\newcommand{\ImportTok}[1]{\textcolor[rgb]{0.00,0.46,0.62}{#1}}
\newcommand{\InformationTok}[1]{\textcolor[rgb]{0.37,0.37,0.37}{#1}}
\newcommand{\KeywordTok}[1]{\textcolor[rgb]{0.00,0.23,0.31}{\textbf{#1}}}
\newcommand{\NormalTok}[1]{\textcolor[rgb]{0.00,0.23,0.31}{#1}}
\newcommand{\OperatorTok}[1]{\textcolor[rgb]{0.37,0.37,0.37}{#1}}
\newcommand{\OtherTok}[1]{\textcolor[rgb]{0.00,0.23,0.31}{#1}}
\newcommand{\PreprocessorTok}[1]{\textcolor[rgb]{0.68,0.00,0.00}{#1}}
\newcommand{\RegionMarkerTok}[1]{\textcolor[rgb]{0.00,0.23,0.31}{#1}}
\newcommand{\SpecialCharTok}[1]{\textcolor[rgb]{0.37,0.37,0.37}{#1}}
\newcommand{\SpecialStringTok}[1]{\textcolor[rgb]{0.13,0.47,0.30}{#1}}
\newcommand{\StringTok}[1]{\textcolor[rgb]{0.13,0.47,0.30}{#1}}
\newcommand{\VariableTok}[1]{\textcolor[rgb]{0.07,0.07,0.07}{#1}}
\newcommand{\VerbatimStringTok}[1]{\textcolor[rgb]{0.13,0.47,0.30}{#1}}
\newcommand{\WarningTok}[1]{\textcolor[rgb]{0.37,0.37,0.37}{\textit{#1}}}

\usepackage{longtable,booktabs,array}
\newcounter{none} % for unnumbered tables
\usepackage{calc} % for calculating minipage widths
% Correct order of tables after \paragraph or \subparagraph
\usepackage{etoolbox}
\makeatletter
\patchcmd\longtable{\par}{\if@noskipsec\mbox{}\fi\par}{}{}
\makeatother
% Allow footnotes in longtable head/foot
\IfFileExists{footnotehyper.sty}{\usepackage{footnotehyper}}{\usepackage{footnote}}
\makesavenoteenv{longtable}
\usepackage{graphicx}
\makeatletter
\newsavebox\pandoc@box
\newcommand*\pandocbounded[1]{% scales image to fit in text height/width
  \sbox\pandoc@box{#1}%
  \Gscale@div\@tempa{\textheight}{\dimexpr\ht\pandoc@box+\dp\pandoc@box\relax}%
  \Gscale@div\@tempb{\linewidth}{\wd\pandoc@box}%
  \ifdim\@tempb\p@<\@tempa\p@\let\@tempa\@tempb\fi% select the smaller of both
  \ifdim\@tempa\p@<\p@\scalebox{\@tempa}{\usebox\pandoc@box}%
  \else\usebox{\pandoc@box}%
  \fi%
}
% Set default figure placement to htbp
\def\fps@figure{htbp}
\makeatother





\setlength{\emergencystretch}{3em} % prevent overfull lines

\providecommand{\tightlist}{%
  \setlength{\itemsep}{0pt}\setlength{\parskip}{0pt}}



 


\KOMAoption{captions}{tableheading}
\usepackage{amsmath}
\usepackage{xcolor}
\definecolor{azulime}{HTML}{1C3B6F}
\usepackage{hyperref}
\AtBeginDocument{\hypersetup{colorlinks=true, urlcolor=azulime,linkcolor=azulime, citecolor=azulime}}
\usepackage{fancyhdr}
\usepackage{graphicx}
\pagestyle{fancy}
\fancyhf{}
\fancyhead[R]{\includegraphics[width=2.5cm]{logo.png}}
\fancyfoot[R]{\thepage}
\renewcommand{\headrulewidth}{0pt}
\renewcommand{\footrulewidth}{0pt}
\usepackage{titling}
\pretitle{\begin{center}\includegraphics[width=4cm]{logo.png}\\[1cm]}
\posttitle{\end{center}\thispagestyle{empty}\setcounter{page}{1}}
\makeatletter
\@ifpackageloaded{caption}{}{\usepackage{caption}}
\AtBeginDocument{%
\ifdefined\contentsname
  \renewcommand*\contentsname{Table of contents}
\else
  \newcommand\contentsname{Table of contents}
\fi
\ifdefined\listfigurename
  \renewcommand*\listfigurename{List of Figures}
\else
  \newcommand\listfigurename{List of Figures}
\fi
\ifdefined\listtablename
  \renewcommand*\listtablename{List of Tables}
\else
  \newcommand\listtablename{List of Tables}
\fi
\ifdefined\figurename
  \renewcommand*\figurename{Figure}
\else
  \newcommand\figurename{Figure}
\fi
\ifdefined\tablename
  \renewcommand*\tablename{Table}
\else
  \newcommand\tablename{Table}
\fi
}
\@ifpackageloaded{float}{}{\usepackage{float}}
\floatstyle{ruled}
\@ifundefined{c@chapter}{\newfloat{codelisting}{h}{lop}}{\newfloat{codelisting}{h}{lop}[chapter]}
\floatname{codelisting}{Listing}
\newcommand*\listoflistings{\listof{codelisting}{List of Listings}}
\makeatother
\makeatletter
\makeatother
\makeatletter
\@ifpackageloaded{caption}{}{\usepackage{caption}}
\@ifpackageloaded{subcaption}{}{\usepackage{subcaption}}
\makeatother
\usepackage{bookmark}
\IfFileExists{xurl.sty}{\usepackage{xurl}}{} % add URL line breaks if available
\urlstyle{same}
\hypersetup{
  pdfauthor={Jaqueline Lamas da Silva (17646475),; Olian Castanha Lima (13686261),; Oyindamola Olubusola Alalade (17647114)},
  colorlinks=true,
  linkcolor={blue},
  filecolor={Maroon},
  citecolor={Blue},
  urlcolor={Blue},
  pdfcreator={LaTeX via pandoc}}


\title{\LARGE \textbf{MAE5834 - Estudo de Simulação Inferência Robusta}}
\author{Jaqueline Lamas da Silva (17646475), \and Olian Castanha Lima
(13686261), \and Oyindamola Olubusola Alalade (17647114)}
\date{}
\begin{document}
\maketitle

\renewcommand*\contentsname{Sumário}
{
\setcounter{tocdepth}{3}
\tableofcontents
}

\newpage

\section{Introdução}\label{introduuxe7uxe3o}

Este trabalho prático tem como objetivo avaliar o desempenho do MDPDE em
comparação ao EMV sob diferentes cenários de contaminação e valores da
constante de afinação \(\alpha \in [0,1)\). O modelo sorteado foi o
Gamma.

\subsubsection{(a)}\label{a}

\textbf{Para o modelo sorteado, escreva a função de densidade, o vetor
de escore e a função objetivo do MDPDE a ser minimizada.}

Seja Y uma variável aleatória com distribuição gama de média \(\mu\) e
coeficiente de variação \(\phi^{-\frac{1}{2}}\). A função de densidade
de probabilidade de Y é dada por:

\begin{align*}
f_Y(y,\mu,\phi) = exp\left\{ \phi \left[ \frac{-y}{\mu} - log(\mu) \right] - log(\Gamma(\phi)) + \phi log(\phi y) - log(y) \right\},
\end{align*} em que \(y>0\), \(\phi>0\), \(\mu>0\) e
\(\Gamma(\phi) = \int_0^\infty t ^{\phi -1} e^{-t}dt\) é a função gama.

\textbf{Vetor Escore}

Seja \(Y_1, \cdots, Y_n\) uma amostra aleatória de Y

\begin{align*}
Y_i|\mathbf{x_i} \overset{\text{ind}}{\sim} Gamma (\mu_i, \phi) \\
\mu_i = log(\eta_i), \eta_i = \mathbf{x_i}^{\intercal} \boldsymbol{\beta}
\end{align*}

Com isso, o vetor escore para a i-ésima observação é dado por:
\begin{align*} 
u_{\boldsymbol{\beta}}(y_i) = \phi \frac{(y_i - \mu_i)}{\mu_i} \mathbf{x_i},
\end{align*} em que \(\mathbf{x_i}\) é a i-ésima linha da matriz
\(\mathbf{X}\).

\textbf{Função objetivo}

O MDPDE minimiza a seguinte função perda:

\begin{align}\label{eq.1}
H_n(\boldsymbol{\beta}) = \frac{1}{n} \sum_{i=1}^n \left[ \int_{\mathcal{X}} f_Y^{1+\alpha}(z; \mu_i, \phi) \, dz - \left( 1 + \frac{1}{\alpha} \right) f_Y^\alpha(Y_i; \mu_i, \phi) \right], \text{ para } \alpha \in [0,1).
\end{align}

Então a função score é: \begin{align*}
\psi(y_i,\beta) = u_{\beta}(y_i) f_Y^{\alpha}(y_i;\mu_i, \phi)-\int_{\mathcal{X}}  u_{\beta}(z) f_Y^{1+\alpha}(z;\mu, \phi)(z)dz
\end{align*}

E as estimativas são obtidas resolvendo o sistema de equações de
estimação: \[\sum_{i=1}^n \psi_i(y_i, \beta) = \mathbf{0}\]

\subsubsection{(b)}\label{b}

Modelo verdadeiro: \begin{align*}
Y_i|\mathbf{x_i} \overset{\text{ind}}{\sim} Gamma (\mu_i, \phi) \\
\mu_i = log(\eta_i), \eta_i = \beta_0 + \beta_1 x_{i1}
\end{align*}, em que

\[\boldsymbol{\beta} = (0.3, 0.1)^\intercal;\ \phi=4;\ x_{i1} \sim U(1,10);\ i=1,\cdots,n. \]

\subsubsection{(c)}\label{c}

Mecanismos de Contaminação: Para cada amostra gerada, considere os
seguintes cenários de contaminação com proporção
\(\epsilon in {0%, 5%, 10%}
\):

\begin{enumerate}
\def\labelenumi{\roman{enumi})}
\tightlist
\item
  Cenário C0 (Sem contaminação): Dados puros gerados estritamente pelo
  modelo teórico.
\end{enumerate}

\begin{Shaded}
\begin{Highlighting}[]
\NormalTok{alfas}\OtherTok{\textless{}{-}}\FunctionTok{c}\NormalTok{(}\DecValTok{0}\NormalTok{,}\FloatTok{0.1}\NormalTok{,}\FloatTok{0.25}\NormalTok{,}\FloatTok{0.5}\NormalTok{,}\FloatTok{0.75}\NormalTok{,}\DecValTok{1}\NormalTok{)}
\NormalTok{nome.amostra}\OtherTok{\textless{}{-}}\StringTok{"C0"}

\FunctionTok{jpeg}\NormalTok{(}\AttributeTok{file=}\FunctionTok{paste0}\NormalTok{(}\StringTok{"Resultados/"}\NormalTok{,}\FunctionTok{paste0}\NormalTok{(nome.amostra,}\StringTok{"\_Intercepto"}\NormalTok{ ,}\StringTok{".png"}\NormalTok{)), }\AttributeTok{width =} \DecValTok{1000}\NormalTok{, }\AttributeTok{height =} \DecValTok{500}\NormalTok{, }\AttributeTok{quality =} \DecValTok{100}\NormalTok{, }\AttributeTok{pointsize =} \DecValTok{20}\NormalTok{)}
\FunctionTok{par}\NormalTok{(}\AttributeTok{mfrow=}\FunctionTok{c}\NormalTok{(}\DecValTok{2}\NormalTok{,}\DecValTok{3}\NormalTok{), }\AttributeTok{mar=}\FunctionTok{c}\NormalTok{(}\DecValTok{3}\NormalTok{,}\DecValTok{3}\NormalTok{,}\DecValTok{3}\NormalTok{, }\DecValTok{2}\NormalTok{))}
\NormalTok{boxplot.intercepto}\OtherTok{\textless{}{-}}\ControlFlowTok{function}\NormalTok{(alpha)}
\NormalTok{\{}
\NormalTok{  recupera.intercepto.emv}\OtherTok{\textless{}{-}}\ControlFlowTok{function}\NormalTok{(n)}
\NormalTok{  \{}
    
\NormalTok{    arquivo}\OtherTok{\textless{}{-}}\FunctionTok{paste0}\NormalTok{(}\StringTok{"Amostras/"}\NormalTok{, nome.amostra, }\StringTok{"\_"}\NormalTok{, n,}\StringTok{"("}\NormalTok{,alpha,}\StringTok{")"}\NormalTok{, }\StringTok{".RDS"}\NormalTok{)}
\NormalTok{    amostras}\OtherTok{\textless{}{-}}\FunctionTok{readRDS}\NormalTok{(arquivo)}
\NormalTok{    interceptos.emv}\OtherTok{\textless{}{-}}\FunctionTok{sapply}\NormalTok{(}\DecValTok{1}\SpecialCharTok{:}\NormalTok{M,}\AttributeTok{FUN=}\ControlFlowTok{function}\NormalTok{(x)\{amostras[[x]]}\SpecialCharTok{$}\NormalTok{theta.emv[}\DecValTok{1}\NormalTok{]\})}
    \FunctionTok{return}\NormalTok{(interceptos.emv)}
\NormalTok{  \}}
\NormalTok{  interceptos.emv}\OtherTok{\textless{}{-}}\FunctionTok{lapply}\NormalTok{(tam.n, recupera.intercepto.emv)}
  
\NormalTok{  recupera.intercepto.emdpde}\OtherTok{\textless{}{-}}\ControlFlowTok{function}\NormalTok{(n)}
\NormalTok{  \{}
    
\NormalTok{    arquivo}\OtherTok{\textless{}{-}}\FunctionTok{paste0}\NormalTok{(}\StringTok{"Amostras/"}\NormalTok{, nome.amostra, }\StringTok{"\_"}\NormalTok{, n,}\StringTok{"("}\NormalTok{,alpha,}\StringTok{")"}\NormalTok{, }\StringTok{".RDS"}\NormalTok{)}
\NormalTok{    amostras}\OtherTok{\textless{}{-}}\FunctionTok{readRDS}\NormalTok{(arquivo)}
\NormalTok{    interceptos.emdpde}\OtherTok{\textless{}{-}}\FunctionTok{sapply}\NormalTok{(}\DecValTok{1}\SpecialCharTok{:}\NormalTok{M,}\AttributeTok{FUN=}\ControlFlowTok{function}\NormalTok{(x)\{amostras[[x]]}\SpecialCharTok{$}\NormalTok{theta.emdpde[}\DecValTok{1}\NormalTok{]\})}
    \FunctionTok{return}\NormalTok{(interceptos.emdpde)}
\NormalTok{  \}}
\NormalTok{  interceptos.emdpde}\OtherTok{\textless{}{-}}\FunctionTok{lapply}\NormalTok{(tam.n, recupera.intercepto.emdpde)}
  
\NormalTok{  paleta\_completa }\OtherTok{\textless{}{-}} \FunctionTok{c}\NormalTok{(}\StringTok{"\#EADDF0"}\NormalTok{, }\StringTok{"\#C2A5CF"}\NormalTok{, }\StringTok{"\#8B629C"}\NormalTok{, }\StringTok{"\#E8F5D6"}\NormalTok{, }\StringTok{"\#C9E69E"}\NormalTok{, }\StringTok{"\#8CB858"}\NormalTok{)}
  \FunctionTok{boxplot}\NormalTok{(}\FunctionTok{c}\NormalTok{(interceptos.emv,interceptos.emdpde),}\AttributeTok{boxwex=}\FloatTok{0.2}\NormalTok{,}\AttributeTok{main=}\FunctionTok{paste0}\NormalTok{(nome.amostra,}\StringTok{" Intercepto"}\NormalTok{ ,}\StringTok{" ("}\NormalTok{, alpha,}\StringTok{")"}\NormalTok{), }\AttributeTok{pch=}\DecValTok{1}\NormalTok{,}\AttributeTok{at=}\FunctionTok{c}\NormalTok{(}\DecValTok{1}\NormalTok{,}\FloatTok{1.25}\NormalTok{,}\FloatTok{1.5}\NormalTok{, }\FloatTok{1.75}\NormalTok{, }\DecValTok{2}\NormalTok{, }\FloatTok{2.25}\NormalTok{), }\AttributeTok{xaxt =} \StringTok{"n"}\NormalTok{,}\AttributeTok{col=}\NormalTok{paleta\_completa)}
  \FunctionTok{abline}\NormalTok{(}\AttributeTok{h=}\NormalTok{beta.verd[}\DecValTok{1}\NormalTok{], }\AttributeTok{lty=}\DecValTok{2}\NormalTok{, }\AttributeTok{lwd=}\DecValTok{2}\NormalTok{)}
  
  \ControlFlowTok{if}\NormalTok{(alpha}\SpecialCharTok{==}\DecValTok{0}\NormalTok{)}\FunctionTok{legend}\NormalTok{(}\StringTok{"topright"}\NormalTok{,}\AttributeTok{legend =} \FunctionTok{paste0}\NormalTok{(}\FunctionTok{rep}\NormalTok{(}\FunctionTok{c}\NormalTok{(}\StringTok{"EMV"}\NormalTok{,}\StringTok{"MDPDE"}\NormalTok{),}\AttributeTok{each=}\DecValTok{3}\NormalTok{),}\StringTok{" {-} n = "}\NormalTok{,tam.n), }\AttributeTok{col=}\NormalTok{paleta\_completa, }\AttributeTok{bty=}\StringTok{"n"}\NormalTok{, }\AttributeTok{fill=}\NormalTok{paleta\_completa, }\AttributeTok{cex=}\FloatTok{0.5}\NormalTok{)}
\NormalTok{\}}
\FunctionTok{lapply}\NormalTok{(alfas, boxplot.intercepto)}
\end{Highlighting}
\end{Shaded}

\begin{verbatim}
[[1]]
[[1]]$rect
[[1]]$rect$w
[1] 0.7611472

[[1]]$rect$h
[1] 0.1522294

[[1]]$rect$left
[1] 2.078853

[[1]]$rect$top
[1] 0.5784051


[[1]]$text
[[1]]$text$x
[1] 2.237886 2.237886 2.237886 2.237886 2.237886 2.237886

[[1]]$text$y
[1] 0.5566580 0.5349109 0.5131639 0.4914168 0.4696698 0.4479227



[[2]]
NULL

[[3]]
NULL

[[4]]
NULL

[[5]]
NULL

[[6]]
NULL
\end{verbatim}

\begin{Shaded}
\begin{Highlighting}[]
\FunctionTok{dev.off}\NormalTok{()}
\end{Highlighting}
\end{Shaded}

\begin{verbatim}
pdf 
  2 
\end{verbatim}

\begin{Shaded}
\begin{Highlighting}[]
\FunctionTok{jpeg}\NormalTok{(}\AttributeTok{file=}\FunctionTok{paste0}\NormalTok{(}\StringTok{"Resultados/"}\NormalTok{,}\FunctionTok{paste0}\NormalTok{(nome.amostra,}\StringTok{"\_Beta1"}\NormalTok{ ,}\StringTok{".png"}\NormalTok{)), }\AttributeTok{width =} \DecValTok{1000}\NormalTok{, }\AttributeTok{height =} \DecValTok{500}\NormalTok{, }\AttributeTok{quality =} \DecValTok{100}\NormalTok{, }\AttributeTok{pointsize =} \DecValTok{20}\NormalTok{)}
\FunctionTok{par}\NormalTok{(}\AttributeTok{mfrow=}\FunctionTok{c}\NormalTok{(}\DecValTok{2}\NormalTok{,}\DecValTok{3}\NormalTok{), }\AttributeTok{mar=}\FunctionTok{c}\NormalTok{(}\DecValTok{3}\NormalTok{,}\DecValTok{3}\NormalTok{,}\DecValTok{3}\NormalTok{, }\DecValTok{2}\NormalTok{))}
\NormalTok{boxplot.beta1}\OtherTok{\textless{}{-}}\ControlFlowTok{function}\NormalTok{(alfa)}
\NormalTok{\{}
\NormalTok{  recupera.beta1.emv}\OtherTok{\textless{}{-}}\ControlFlowTok{function}\NormalTok{(n)}
\NormalTok{  \{}
    
\NormalTok{    arquivo}\OtherTok{\textless{}{-}}\FunctionTok{paste0}\NormalTok{(}\StringTok{"Amostras/"}\NormalTok{, nome.amostra, }\StringTok{"\_"}\NormalTok{, n,}\StringTok{"("}\NormalTok{,alfa,}\StringTok{")"}\NormalTok{, }\StringTok{".RDS"}\NormalTok{)}
\NormalTok{    amostras}\OtherTok{\textless{}{-}}\FunctionTok{readRDS}\NormalTok{(arquivo)}
\NormalTok{    betas.emv}\OtherTok{\textless{}{-}}\FunctionTok{sapply}\NormalTok{(}\DecValTok{1}\SpecialCharTok{:}\NormalTok{M,}\AttributeTok{FUN=}\ControlFlowTok{function}\NormalTok{(x)\{amostras[[x]]}\SpecialCharTok{$}\NormalTok{theta.emv[}\DecValTok{1}\NormalTok{]\})}
    \FunctionTok{return}\NormalTok{(betas.emv)}
\NormalTok{  \}}
\NormalTok{  betas.emv}\OtherTok{\textless{}{-}}\FunctionTok{lapply}\NormalTok{(tam.n, recupera.beta1.emv)}
  
\NormalTok{  recupera.beta1.emdpde}\OtherTok{\textless{}{-}}\ControlFlowTok{function}\NormalTok{(n)}
\NormalTok{  \{}
    
\NormalTok{    arquivo}\OtherTok{\textless{}{-}}\FunctionTok{paste0}\NormalTok{(}\StringTok{"Amostras/"}\NormalTok{, nome.amostra, }\StringTok{"\_"}\NormalTok{, n,}\StringTok{"("}\NormalTok{,alfa,}\StringTok{")"}\NormalTok{, }\StringTok{".RDS"}\NormalTok{)}
\NormalTok{    amostras}\OtherTok{\textless{}{-}}\FunctionTok{readRDS}\NormalTok{(arquivo)}
\NormalTok{    betas.emdpde}\OtherTok{\textless{}{-}}\FunctionTok{sapply}\NormalTok{(}\DecValTok{1}\SpecialCharTok{:}\NormalTok{M,}\AttributeTok{FUN=}\ControlFlowTok{function}\NormalTok{(x)\{amostras[[x]]}\SpecialCharTok{$}\NormalTok{theta.emdpde[}\DecValTok{1}\NormalTok{]\})}
    \FunctionTok{return}\NormalTok{(betas.emdpde)}
\NormalTok{  \}}
\NormalTok{  betas.emdpde}\OtherTok{\textless{}{-}}\FunctionTok{lapply}\NormalTok{(tam.n, recupera.beta1.emdpde)}
  
\NormalTok{  paleta\_completa }\OtherTok{\textless{}{-}} \FunctionTok{c}\NormalTok{(}\StringTok{"\#EADDF0"}\NormalTok{, }\StringTok{"\#C2A5CF"}\NormalTok{, }\StringTok{"\#8B629C"}\NormalTok{, }\StringTok{"\#E8F5D6"}\NormalTok{, }\StringTok{"\#C9E69E"}\NormalTok{, }\StringTok{"\#8CB858"}\NormalTok{)}
  \FunctionTok{boxplot}\NormalTok{(}\FunctionTok{c}\NormalTok{(betas.emv,betas.emdpde),}\AttributeTok{boxwex=}\FloatTok{0.2}\NormalTok{,}\AttributeTok{main=}\FunctionTok{paste0}\NormalTok{(nome.amostra,}\StringTok{" Beta1"}\NormalTok{ ,}\StringTok{" ("}\NormalTok{, alfa,}\StringTok{")"}\NormalTok{), }\AttributeTok{pch=}\DecValTok{1}\NormalTok{,}\AttributeTok{at=}\FunctionTok{c}\NormalTok{(}\DecValTok{1}\NormalTok{,}\FloatTok{1.25}\NormalTok{,}\FloatTok{1.5}\NormalTok{, }\FloatTok{1.75}\NormalTok{, }\DecValTok{2}\NormalTok{, }\FloatTok{2.25}\NormalTok{), }\AttributeTok{xaxt =} \StringTok{"n"}\NormalTok{,}\AttributeTok{col=}\NormalTok{paleta\_completa)}
  \FunctionTok{abline}\NormalTok{(}\AttributeTok{h=}\NormalTok{beta.verd[}\DecValTok{1}\NormalTok{], }\AttributeTok{lty=}\DecValTok{2}\NormalTok{, }\AttributeTok{lwd=}\DecValTok{2}\NormalTok{)}
  
  \ControlFlowTok{if}\NormalTok{(alfa}\SpecialCharTok{==}\DecValTok{0}\NormalTok{)}\FunctionTok{legend}\NormalTok{(}\StringTok{"topright"}\NormalTok{,}\AttributeTok{legend =} \FunctionTok{paste0}\NormalTok{(}\FunctionTok{rep}\NormalTok{(}\FunctionTok{c}\NormalTok{(}\StringTok{"EMV"}\NormalTok{,}\StringTok{"MDPDE"}\NormalTok{),}\AttributeTok{each=}\DecValTok{3}\NormalTok{),}\StringTok{" {-} n = "}\NormalTok{,tam.n), }\AttributeTok{col=}\NormalTok{paleta\_completa, }\AttributeTok{bty=}\StringTok{"n"}\NormalTok{, }\AttributeTok{fill=}\NormalTok{paleta\_completa, }\AttributeTok{cex=}\FloatTok{0.5}\NormalTok{)}
\NormalTok{\}}
\FunctionTok{lapply}\NormalTok{(alfas, boxplot.beta1)}
\end{Highlighting}
\end{Shaded}

\begin{verbatim}
[[1]]
[[1]]$rect
[[1]]$rect$w
[1] 0.7611472

[[1]]$rect$h
[1] 0.1522294

[[1]]$rect$left
[1] 2.078853

[[1]]$rect$top
[1] 0.5784051


[[1]]$text
[[1]]$text$x
[1] 2.237886 2.237886 2.237886 2.237886 2.237886 2.237886

[[1]]$text$y
[1] 0.5566580 0.5349109 0.5131639 0.4914168 0.4696698 0.4479227



[[2]]
NULL

[[3]]
NULL

[[4]]
NULL

[[5]]
NULL

[[6]]
NULL
\end{verbatim}

\begin{Shaded}
\begin{Highlighting}[]
\FunctionTok{dev.off}\NormalTok{()}
\end{Highlighting}
\end{Shaded}

\begin{verbatim}
pdf 
  2 
\end{verbatim}

\begin{enumerate}
\def\labelenumi{\roman{enumi})}
\setcounter{enumi}{1}
\tightlist
\item
  Cenário C1 (Outliers verticais na resposta Y): Uma fração \(\epsilon\)
  das observações é substituída por valores atípicos severos na resposta
  Y.
\end{enumerate}

Neste cenário uma fração das observações tiveram suas respostas
substituídas por um valor gerado de uma gamma com média 5 vezes o valor
original de média para a observação.

\subsubsection{5\% de contaminação na
resposta}\label{de-contaminauxe7uxe3o-na-resposta}

\begin{Shaded}
\begin{Highlighting}[]
\NormalTok{nome.amostra}\OtherTok{\textless{}{-}}\StringTok{"C1.5"}

\FunctionTok{jpeg}\NormalTok{(}\AttributeTok{file=}\FunctionTok{paste0}\NormalTok{(}\StringTok{"Resultados/"}\NormalTok{,}\FunctionTok{paste0}\NormalTok{(nome.amostra,}\StringTok{"\_Intercepto"}\NormalTok{ ,}\StringTok{".png"}\NormalTok{)), }\AttributeTok{width =} \DecValTok{1000}\NormalTok{, }\AttributeTok{height =} \DecValTok{500}\NormalTok{, }\AttributeTok{quality =} \DecValTok{100}\NormalTok{, }\AttributeTok{pointsize =} \DecValTok{20}\NormalTok{)}
\FunctionTok{par}\NormalTok{(}\AttributeTok{mfrow=}\FunctionTok{c}\NormalTok{(}\DecValTok{2}\NormalTok{,}\DecValTok{3}\NormalTok{), }\AttributeTok{mar=}\FunctionTok{c}\NormalTok{(}\DecValTok{3}\NormalTok{,}\DecValTok{3}\NormalTok{,}\DecValTok{3}\NormalTok{, }\DecValTok{2}\NormalTok{))}
\FunctionTok{lapply}\NormalTok{(alfas, boxplot.intercepto)}
\end{Highlighting}
\end{Shaded}

\begin{verbatim}
[[1]]
[[1]]$rect
[[1]]$rect$w
[1] 0.7611472

[[1]]$rect$h
[1] 0.2397439

[[1]]$rect$left
[1] 2.078853

[[1]]$rect$top
[1] 0.697799


[[1]]$text
[[1]]$text$x
[1] 2.237886 2.237886 2.237886 2.237886 2.237886 2.237886

[[1]]$text$y
[1] 0.6635499 0.6293007 0.5950516 0.5608025 0.5265533 0.4923042



[[2]]
NULL

[[3]]
NULL

[[4]]
NULL

[[5]]
NULL

[[6]]
NULL
\end{verbatim}

\begin{Shaded}
\begin{Highlighting}[]
\FunctionTok{dev.off}\NormalTok{()}
\end{Highlighting}
\end{Shaded}

\begin{verbatim}
pdf 
  2 
\end{verbatim}

\begin{Shaded}
\begin{Highlighting}[]
\FunctionTok{jpeg}\NormalTok{(}\AttributeTok{file=}\FunctionTok{paste0}\NormalTok{(}\StringTok{"Resultados/"}\NormalTok{,}\FunctionTok{paste0}\NormalTok{(nome.amostra,}\StringTok{"\_Beta1"}\NormalTok{ ,}\StringTok{".png"}\NormalTok{)), }\AttributeTok{width =} \DecValTok{1000}\NormalTok{, }\AttributeTok{height =} \DecValTok{500}\NormalTok{, }\AttributeTok{quality =} \DecValTok{100}\NormalTok{, }\AttributeTok{pointsize =} \DecValTok{20}\NormalTok{)}
\FunctionTok{par}\NormalTok{(}\AttributeTok{mfrow=}\FunctionTok{c}\NormalTok{(}\DecValTok{2}\NormalTok{,}\DecValTok{3}\NormalTok{), }\AttributeTok{mar=}\FunctionTok{c}\NormalTok{(}\DecValTok{3}\NormalTok{,}\DecValTok{3}\NormalTok{,}\DecValTok{3}\NormalTok{, }\DecValTok{2}\NormalTok{))}
\FunctionTok{lapply}\NormalTok{(alfas, boxplot.beta1)}
\end{Highlighting}
\end{Shaded}

\begin{verbatim}
[[1]]
[[1]]$rect
[[1]]$rect$w
[1] 0.7611472

[[1]]$rect$h
[1] 0.2397439

[[1]]$rect$left
[1] 2.078853

[[1]]$rect$top
[1] 0.697799


[[1]]$text
[[1]]$text$x
[1] 2.237886 2.237886 2.237886 2.237886 2.237886 2.237886

[[1]]$text$y
[1] 0.6635499 0.6293007 0.5950516 0.5608025 0.5265533 0.4923042



[[2]]
NULL

[[3]]
NULL

[[4]]
NULL

[[5]]
NULL

[[6]]
NULL
\end{verbatim}

\begin{Shaded}
\begin{Highlighting}[]
\FunctionTok{dev.off}\NormalTok{()}
\end{Highlighting}
\end{Shaded}

\begin{verbatim}
pdf 
  2 
\end{verbatim}

\subsubsection{10\% de contaminação na
resposta}\label{de-contaminauxe7uxe3o-na-resposta-1}

\begin{Shaded}
\begin{Highlighting}[]
\NormalTok{nome.amostra}\OtherTok{\textless{}{-}}\StringTok{"C1.10"}

\FunctionTok{jpeg}\NormalTok{(}\AttributeTok{file=}\FunctionTok{paste0}\NormalTok{(}\StringTok{"Resultados/"}\NormalTok{,}\FunctionTok{paste0}\NormalTok{(nome.amostra,}\StringTok{"\_Intercepto"}\NormalTok{ ,}\StringTok{".png"}\NormalTok{)), }\AttributeTok{width =} \DecValTok{1000}\NormalTok{, }\AttributeTok{height =} \DecValTok{500}\NormalTok{, }\AttributeTok{quality =} \DecValTok{100}\NormalTok{, }\AttributeTok{pointsize =} \DecValTok{20}\NormalTok{)}
\FunctionTok{par}\NormalTok{(}\AttributeTok{mfrow=}\FunctionTok{c}\NormalTok{(}\DecValTok{2}\NormalTok{,}\DecValTok{3}\NormalTok{), }\AttributeTok{mar=}\FunctionTok{c}\NormalTok{(}\DecValTok{3}\NormalTok{,}\DecValTok{3}\NormalTok{,}\DecValTok{3}\NormalTok{, }\DecValTok{2}\NormalTok{))}
\FunctionTok{lapply}\NormalTok{(alfas, boxplot.intercepto)}
\FunctionTok{dev.off}\NormalTok{()}

\FunctionTok{jpeg}\NormalTok{(}\AttributeTok{file=}\FunctionTok{paste0}\NormalTok{(}\StringTok{"Resultados/"}\NormalTok{,}\FunctionTok{paste0}\NormalTok{(nome.amostra,}\StringTok{"\_Beta1"}\NormalTok{ ,}\StringTok{".png"}\NormalTok{)), }\AttributeTok{width =} \DecValTok{1000}\NormalTok{, }\AttributeTok{height =} \DecValTok{500}\NormalTok{, }\AttributeTok{quality =} \DecValTok{100}\NormalTok{, }\AttributeTok{pointsize =} \DecValTok{20}\NormalTok{)}
\FunctionTok{par}\NormalTok{(}\AttributeTok{mfrow=}\FunctionTok{c}\NormalTok{(}\DecValTok{2}\NormalTok{,}\DecValTok{3}\NormalTok{), }\AttributeTok{mar=}\FunctionTok{c}\NormalTok{(}\DecValTok{3}\NormalTok{,}\DecValTok{3}\NormalTok{,}\DecValTok{3}\NormalTok{, }\DecValTok{2}\NormalTok{))}
\FunctionTok{lapply}\NormalTok{(alfas, boxplot.beta1)}
\FunctionTok{dev.off}\NormalTok{()}
\end{Highlighting}
\end{Shaded}

\begin{enumerate}
\def\labelenumi{\roman{enumi})}
\setcounter{enumi}{2}
\item
  Cenário C2 (Pontos de alavancagem em X): Uma fração \(\epsilon\) do
  vetor de covariáveis contínuas é substituída por pontos
  extremos/distantes. Neste cenário uma fração das observações tiveram
  seus respectivos valores para a covariável \(x_{i1}\) gerada de uma
  uniforme no intervalo (16,17).
\item
  Cenário C3 (Contaminação mista X e Y): Combinação simultânea de pontos
  de alavancagem em xi1 com resposta discrepante em Y. Neste cenário uma
  fração das observação tiveram simultaneamente suas repostas
  substituídas por um valor gerado de uma gamma com média 5 vezes o
  valor original e seus respectivos valores para a covariável \(x_{i1}\)
  gerados de uma uniforme no intervalo (16,17).
\end{enumerate}

\section{Referências}\label{referuxeancias}




\end{document}
